\documentclass[10pt,a4paper]{amsart}
\usepackage{amsmath,amssymb,amsthm,graphicx}
% Retained minimal preamble for the Stacks Project chapter coherent.tex (Stacks chapter, pre-delegation retained fragment for stacks-coherent).
% Verbatim excerpt of the chapter preamble (sources/stacks-spaces-descent/preamble.tex): the reference/slogan/history
% environment block (native lines 7-12), Xy-pic setup (native lines 14-19) and the native theorem declarations with
% the mathematical macros (native lines 154-217). Copied mechanically; nothing added, removed or redefined. The
% wrapper class and the amsmath/amssymb/amsthm/graphicx base packages are supplied by the assembler, so the
% class/xr-hyper/hyperref/externaldocument part of the native preamble is omitted.
% For dealing with references we use the comment environment
\usepackage{verbatim}
\newenvironment{reference}{\comment}{\endcomment}
%\newenvironment{reference}{}{}
\newenvironment{slogan}{\comment}{\endcomment}
\newenvironment{history}{\comment}{\endcomment}
% For commutative diagrams we use Xy-pic
\usepackage[all]{xy}

% We use 2cell for 2-commutative diagrams.
\xyoption{2cell}
\UseAllTwocells
% Theorem environments.
%
\theoremstyle{plain}
\newtheorem{theorem}[subsection]{Theorem}
\newtheorem{proposition}[subsection]{Proposition}
\newtheorem{lemma}[subsection]{Lemma}

\theoremstyle{definition}
\newtheorem{definition}[subsection]{Definition}
\newtheorem{example}[subsection]{Example}
\newtheorem{exercise}[subsection]{Exercise}
\newtheorem{situation}[subsection]{Situation}

\theoremstyle{remark}
\newtheorem{remark}[subsection]{Remark}
\newtheorem{remarks}[subsection]{Remarks}

\numberwithin{equation}{subsection}

% Macros
%
\def\lim{\mathop{\mathrm{lim}}\nolimits}
\def\colim{\mathop{\mathrm{colim}}\nolimits}
\def\Spec{\mathop{\mathrm{Spec}}}
\def\Hom{\mathop{\mathrm{Hom}}\nolimits}
\def\Ext{\mathop{\mathrm{Ext}}\nolimits}
\def\SheafHom{\mathop{\mathcal{H}\!\mathit{om}}\nolimits}
\def\SheafExt{\mathop{\mathcal{E}\!\mathit{xt}}\nolimits}
\def\Sch{\mathit{Sch}}
\def\Mor{\mathop{\mathrm{Mor}}\nolimits}
\def\Ob{\mathop{\mathrm{Ob}}\nolimits}
\def\Sh{\mathop{\mathit{Sh}}\nolimits}
\def\NL{\mathop{N\!L}\nolimits}
\def\CH{\mathop{\mathrm{CH}}\nolimits}
\def\proetale{{pro\text{-}\acute{e}tale}}
\def\etale{{\acute{e}tale}}
\def\QCoh{\mathit{QCoh}}
\def\Ker{\mathop{\mathrm{Ker}}}
\def\Im{\mathop{\mathrm{Im}}}
\def\Coker{\mathop{\mathrm{Coker}}}
\def\Coim{\mathop{\mathrm{Coim}}}

% Boxtimes
%
\DeclareMathSymbol{\boxtimes}{\mathbin}{AMSa}{"02}

%
% Macros for moduli stacks/spaces
%
\def\QCohstack{\mathcal{QC}\!\mathit{oh}}
\def\Cohstack{\mathcal{C}\!\mathit{oh}}
\def\Spacesstack{\mathcal{S}\!\mathit{paces}}
\def\Quotfunctor{\mathrm{Quot}}
\def\Hilbfunctor{\mathrm{Hilb}}
\def\Curvesstack{\mathcal{C}\!\mathit{urves}}
\def\Polarizedstack{\mathcal{P}\!\mathit{olarized}}
\def\Complexesstack{\mathcal{C}\!\mathit{omplexes}}
% \Pic is the operator that assigns to X its picard group, usage \Pic(X)
% \Picardstack_{X/B} denotes the Picard stack of X over B
% \Picardfunctor_{X/B} denotes the Picard functor of X over B
\def\Pic{\mathop{\mathrm{Pic}}\nolimits}
\def\Picardstack{\mathcal{P}\!\mathit{ic}}
\def\Picardfunctor{\mathrm{Pic}}
\def\Deformationcategory{\mathcal{D}\!\mathit{ef}}
\usepackage{cleveref}
\crefname{originallemma}{original lemma}{original lemmas}
\Crefname{originallemma}{Original lemma}{Original lemmas}
\makeatletter
\newlabel{cohomology-lemma-functoriality-cech}{{01FD}{source}{}{sourcepointer-1}{}}
\newlabel{cohomology-lemma-functoriality-cech@cref}{{[originallemma][0][]01FD}{[1][0][]source}{}{}{}}
\makeatother
\crefname{originallemma}{original lemma}{original lemmas}
\Crefname{originallemma}{Original lemma}{Original lemmas}
\makeatletter
\newlabel{equation-identify}{{01XU}{source}{}{sourcepointer-2}{}}
\newlabel{equation-identify@cref}{{[originallemma][0][]01XU}{[1][0][]source}{}{}{}}
\makeatother
\crefname{originallemma}{original lemma}{original lemmas}
\Crefname{originallemma}{Original lemma}{Original lemmas}
\makeatletter
\newlabel{lemma-cohomology-projective-space-over-ring}{{01XT}{source}{}{sourcepointer-3}{}}
\newlabel{lemma-cohomology-projective-space-over-ring@cref}{{[originallemma][0][]01XT}{[1][0][]source}{}{}{}}
\makeatother
\title{Functoriality of the identification of the cohomology of O\_\{P\textasciicircum{}n\_R\}(d): the identifications of Equation (identify) are compatible with base change along ring maps R -> R', and for f in R[T\_0,...,T\_n] homogeneous of degree m the map on cohomology induced by multiplication by f : O(d) -> O(d+m) is multiplication by f on H\textasciicircum{}0 and the contragredient of multiplication by f from (R[T])\_\{-n-1-(d+m)\} to (R[T])\_\{-n-1-d\} on H\textasciicircum{}n.}
\author{The Stacks Project Authors}
\date{stacks-project:coherent.tex:local-snapshot-20261012; GNU FDL 1.2 or later}
\begin{document}
\maketitle
\noindent\textit{Editorial note.} Statement span is lines 2101-2121 of sources/stacks-coherent/pinned-original-source.tex; proof span is lines 2123-2212, the complete author proof (88 body lines).\par\medskip
\noindent\textit{Editorial note.} Navigation stubs: 3 cited-but-not-included labels are supplied as non-load-bearing pointers with their pinned tags where the snapshot tag index has them.\par\medskip
\noindent\textit{Editorial note.} extra\_packages is empty on purpose: the retained preamble already loads xy with options and a second declaration clashes in the assembler wrapper.\par\medskip
\noindent\textit{Editorial note.} Text is retained verbatim, including the author's own typos and placeholders.\par\medskip
\noindent\textit{Editorial note.} Full derivation in span; the parenthetical basis-element verification is written by the author, not omitted.\par\medskip
\noindent\textit{Editorial note.} Span computed from the \textbackslash{}label marker through the enclosing \textbackslash{}end\{lemma\}; precard\_gate RESULT=OK (probe L9999).\par\medskip
\noindent\textit{Allowed standard prerequisites (editorial).}\begin{itemize}
\item Explicit prerequisite boundary. See why-field; every pointer named here occurs verbatim in the retained statement/proof text.
\end{itemize}\medskip
\ifdefined\hypertarget\hypertarget{sourcepointer-1}{}\fi
\noindent\textit{Source reference.} Original lemma 01FD: sources/stacks-coherent/pinned-original-source.tex, source lines 2212--2212. Cited by the retained statement or proof but not part of the collected spans; supplied only as a navigation pointer. The argument does not depend on its content.\par\medskip
\ifdefined\hypertarget\hypertarget{sourcepointer-2}{}\fi
\noindent\textit{Source reference.} Original lemma 01XU: sources/stacks-coherent/pinned-original-source.tex, source lines 2089--2097. Cited by the retained statement or proof but not part of the collected spans; supplied only as a navigation pointer. The argument does not depend on its content.\par\medskip
\ifdefined\hypertarget\hypertarget{sourcepointer-3}{}\fi
\noindent\textit{Source reference.} Original lemma 01XT: sources/stacks-coherent/pinned-original-source.tex, source lines 1820--1823. Cited by the retained statement or proof but not part of the collected spans; supplied only as a navigation pointer. The argument does not depend on its content.\par\medskip
\begin{lemma}
\label{lemma-identify-functorially}
The identifications of Equation (\ref{equation-identify}) are
compatible with base change w.r.t.\ ring maps $R \to R'$.
Moreover, for any $f \in R[T_0, \ldots, T_n]$ homogeneous
of degree $m$ the map multiplication by $f$
$$
\mathcal{O}_{\mathbf{P}^n_R}(d)
\longrightarrow
\mathcal{O}_{\mathbf{P}^n_R}(d + m)
$$
induces the map on the cohomology group via the identifications
of Equation (\ref{equation-identify}) which is multiplication by
$f$ for $H^0$ and the contragredient of multiplication by $f$
$$
(R[T_0, \ldots, T_n])_{-n - 1 - (d + m)}
\longrightarrow
(R[T_0, \ldots, T_n])_{-n - 1 - d}
$$
on $H^n$.
\end{lemma}

\begin{proof}
Suppose that $R \to R'$ is a ring map.
Let $\mathcal{U}$ be the standard affine open covering of $\mathbf{P}^n_R$,
and let $\mathcal{U}'$ be the standard affine open covering of
$\mathbf{P}^n_{R'}$. Note that $\mathcal{U}'$ is the pullback of the covering
$\mathcal{U}$ under the canonical morphism
$\mathbf{P}^n_{R'} \to \mathbf{P}^n_R$. Hence there
is a map of {\v C}ech complexes
$$
\gamma :
\check{\mathcal{C}}_{ord}^\bullet(\mathcal{U},
\mathcal{O}_{\mathbf{P}_R}(d))
\longrightarrow
\check{\mathcal{C}}_{ord}^\bullet(\mathcal{U}',
\mathcal{O}_{\mathbf{P}_{R'}}(d))
$$
which is compatible with the map on cohomology by
Cohomology, Lemma \ref{cohomology-lemma-functoriality-cech}.
It is clear from the computations in the proof of
Lemma \ref{lemma-cohomology-projective-space-over-ring}
that this map of {\v C}ech complexes is compatible with the identifications
of the cohomology groups in question. (Namely the basis elements for
the {\v C}ech complex over $R$ simply map to the corresponding basis elements
for the {\v C}ech complex over $R'$.) Whence the first statement of the lemma.

\medskip\noindent
Now fix the ring $R$ and consider two homogeneous polynomials
$f, g \in R[T_0, \ldots, T_n]$ both of the same degree $m$.
Since cohomology is an additive functor, it is clear that the
map induced by multiplication by $f + g$ is the same as the sum
of the maps induced by multiplication by $f$ and the map induced
by multiplication by $g$. Moreover, since cohomology is a functor,
a similar result holds for multiplication by a product $fg$ where
$f, g$ are both homogeneous (but not necessarily of the same degree).
Hence to verify the second statement of the lemma it suffices to
prove this when $f = x \in R$ or when $f = T_i$.
In the case of multiplication by an element $x \in R$ the result
follows since every cohomology groups or complex in sight has the
structure of an $R$-module or complex of $R$-modules.
Finally, we consider the case of multiplication by $T_i$
as a $\mathcal{O}_{\mathbf{P}^n_R}$-linear map
$$
\mathcal{O}_{\mathbf{P}^n_R}(d)
\longrightarrow
\mathcal{O}_{\mathbf{P}^n_R}(d + 1)
$$
The statement on $H^0$ is clear. For the statement on $H^n$
consider multiplication by $T_i$ as a map on {\v C}ech complexes
$$
\check{\mathcal{C}}_{ord}^\bullet(\mathcal{U},
\mathcal{O}_{\mathbf{P}_R}(d))
\longrightarrow
\check{\mathcal{C}}_{ord}^\bullet(\mathcal{U},
\mathcal{O}_{\mathbf{P}_R}(d + 1))
$$
We are going to use the notation introduced in the proof of
Lemma \ref{lemma-cohomology-projective-space-over-ring}.
We consider the effect of multiplication by $T_i$
in terms of the decompositions
$$
\check{\mathcal{C}}_{ord}^\bullet(\mathcal{U}, \mathcal{O}_{\mathbf{P}_R}(d))
=
\bigoplus\nolimits_{\vec{e} \in \mathbf{Z}^{n + 1}, \ \sum e_i = d}
\check{\mathcal{C}}^\bullet(\vec{e})
$$
and
$$
\check{\mathcal{C}}_{ord}^\bullet(\mathcal{U},
\mathcal{O}_{\mathbf{P}_R}(d + 1))
=
\bigoplus\nolimits_{\vec{e} \in \mathbf{Z}^{n + 1}, \ \sum e_i = d + 1}
\check{\mathcal{C}}^\bullet(\vec{e})
$$
It is clear that it maps the subcomplex
$\check{\mathcal{C}}^\bullet(\vec{e})$ to the subcomplex
$\check{\mathcal{C}}^\bullet(\vec{e} + \vec{b}_i)$ where
$\vec{b}_i = (0, \ldots, 0, 1, 0, \ldots, 0))$ the $i$th basis vector.
In other words, it maps the summand of $H^n$ corresponding to
$\vec{e}$ with $e_i < 0$ and $\sum e_i = d$
to the summand of $H^n$ corresponding to
$\vec{e} + \vec{b}_i$ (which is zero if $e_i + b_i \geq 0$).
It is easy to see that this corresponds exactly to the action
of the contragredient of multiplication by $T_i$ as a map
$$
(R[T_0, \ldots, T_n])_{-n - 1 - (d + 1)}
\longrightarrow
(R[T_0, \ldots, T_n])_{-n - 1 - d}
$$
This proves the lemma.
\end{proof}

\end{document}
